% C+ conclusion — characterization frame, policy implications as illustrative.
% JLEO-Cplus.

\section{Conclusion}
\label{sec:conclusion}

We characterized loser-side concentration as a theoretically grounded
participation-only construct in Brazilian electronic procurement.
The construct is an empirical regularity---$\valFL$ firms in BEC
participate in $\valThresholdStat$ or more tender-items between
$\valYearStart$ and $\valYearEnd$ without ever winning---rationalized
by a separating-equilibrium framework where cover bidders avoid wins
to preserve their type signal and accumulate participation by
deployment frequency rather than by competitive selection. Six
independent lines of evidence support the framework: a
$6\%$ within-item conditional price premium, a $\valPBUgradientRatio\!\!\times$
oversight gradient consistent with the framework's monitoring-cost
prediction, structural Regime~2 selection in both modal regimes,
Bayesian latent-class agreement between FL classification and
bid-distribution-based cover-class assignment ($\valLCAprobCoverFL$ vs
$\valLCAprobCoverNonFL$, a $\valLCAratio\!\!\times$ ratio), repeated
FL-winner pair counts $\valDyadicRatio\!\!\times$ above stratified-permutation
null, and within-stratum CADE-cobidder discrimination at AUC
$\valAUCFLfirm$. None of these lines, alone, identifies a causal
mechanism; jointly, they constrain the alternative readings of the
phenomenon to a narrow class.

\paragraph{Policy implications, illustrative.}
The construct's value is diagnostic, not causal. As an illustrative
bound---conditional on the price gap reflecting the cover-bidding
mechanism the framework specifies, which our matching diagnostics
\emph{do not establish}---counterfactual specifications in
Appendix~\ref{app:extensions} yield welfare gains of
$\valWelfareCrossfit$--$\valWelfareIV$ (approximately $\valWelfareLowPct$--$\valWelfareHighPct$
of total BEC spending) under removal of the minimum-bidder rule.
Three counterfactual designs (CF1: rule removal; CF2: optimal
threshold under $\valBidCostHigh$ investigation cost; CF3:
$\valPctTenRef$ increase in detection probability $\theta_k$) yield
welfare gains in the same order of magnitude. We emphasize that these
are illustrative bounds, not policy estimates; an oversight body
deploying the construct as a triage tool should treat the price gap
as conditional and the welfare numbers as scoping the magnitude of
the question, not its answer.

\paragraph{Administrative pathway.}
The construct's primary contribution is administrative. We propose a
three-stage pathway---Screen $\to$ Triage $\to$ Forensic---in which
the loser-side construct narrows the case-load of oversight bodies
before forensic bid-microdata methods are deployed where data
permit:
\emph{(Screen)} computes $\log(1+\text{tenders\_count})$ on
contract-award records (the binary $\valFL$ rule is the deployable
information-coarsening); flags top-$k$ environments for
priority. Cost: roughly one CPU-second per state-year on consumer
hardware (Section~\ref{sec:robustness}, Operational metrics).
\emph{(Triage)} cross-references flagged firms with adjudication
records (CADE, CGU, TCE) and operational priors. Cost: minutes per
flag. \emph{(Forensic)} deploys bid-distribution methods
(\citet{imhof2019detecting,wallimann2023machine}) where bid
microdata are available. Cost: days per case. The construct's
operational value is in narrowing the population that reaches stage
3, not in replacing forensic analysis.

\paragraph{Generalization.}
The construct is most likely portable to procurement environments
sharing three structural features: commodity-heavy item composition
(allowing repeated participation), electronic platforms with publicly
auditable participant lists, and institutional asymmetry between
winner and loser sides of the bidding pool (e.g., minimum-bidder
rules in convite-style modalities, or analogous quorum requirements
in other jurisdictions). The construct is least likely portable to
heterogeneous works contracts, private-sector procurement without
participant disclosure, and jurisdictions where bid microdata are
routinely audited. Generalization to ComprasNet, other Brazilian
states, or OECD platforms requires external replication with the
same scope discipline and the same six-line evidence stack.

\paragraph{What we do not claim.}
The construct does not detect cartelist firms in general; the AUC
$\valAUCdirectStd$ against direct CADE defendants in the BEC universe
defines the detection scope. The price gap is not an overcharge
estimate; matching diagnostics show it does not survive
overlap-strict reweighting. The framework is an organizing device for
interpreting the empirical patterns; it is not the source of causal
identification. Welfare numbers are illustrative bounds, not policy
estimates. Mechanism heterogeneity is reported but not identified.
The construct is a flag, not a verdict.
